% ============================================================
%  N.1.4  Metrics
% ============================================================
\subsection{Metrics}
\label{sec:metrics}

Four metrics are used throughout the experiments to evaluate each
mapping and architecture configuration.

\paragraph{Energy ($E$).}
The total energy is the sum of all memory-access energies across every
level of the hierarchy, plus the compute energy of all multiply-accumulate
operations:
%
\begin{equation}
  E = \sum_{\ell \in \text{levels}}
      \bigl(e^{\text{rd}}_\ell \cdot R_\ell
          + e^{\text{wr}}_\ell \cdot W_\ell\bigr)
    + \sum_{k \in \text{layers}} e_{\text{mac}} \cdot N_k,
  \label{eq:energy}
\end{equation}
%
where $e^{\text{rd}}_\ell$ and $e^{\text{wr}}_\ell$ are the per-access
read and write energies of level~$\ell$ (obtained from Accelergy
models), $R_\ell$ and $W_\ell$ are the total read and write counts,
$e_{\text{mac}}$ is the energy of a single MAC, and $N_k$ the number
of MACs in layer~$k$.
Energy is reported in micro-joules ($\mu$J).

\paragraph{Latency ($L$).}
Latency is computed as a bottleneck model: each memory level determines
its own latency based on the ratio between the data volume it must
transfer and its available bandwidth, and the overall latency is the
maximum across all levels (including the compute level):
%
\begin{equation}
  L = \max_{\ell \in \text{levels}} L_\ell.
  \label{eq:latency}
\end{equation}
%
When multiple fused layers are executed sequentially, per-layer
latencies are summed within each level before taking the maximum.
Latency is reported in clock cycles (cc).

\paragraph{Energy-Delay Product (EDP).}
The EDP combines energy and latency into a single figure of merit that
penalises both high energy and high latency:
%
\begin{equation}
  \text{EDP} = E \times L,
  \label{eq:edp}
\end{equation}
%
reported in J$\cdot$cc (energy converted to joules, latency in clock
cycles).
When comparing aggregated results (e.g.\ $\Sigma$\,partial segments),
the EDP is computed from the \emph{summed} energy and \emph{summed}
latency as in Equation~\eqref{eq:sum-partial}, not as the sum of
individual per-segment EDPs.

\paragraph{DRAM Traffic.}
DRAM reads and writes count the number of operand accesses transferred
between the accelerator and external memory.
This metric directly reflects the benefit of layer fusion:
in single-layer execution the output activations of layer~$i$ are
written to DRAM and then read back as inputs for layer~$i+1$,
whereas in fused execution intermediate activations are kept
entirely on-chip (DRAM bypasses \texttt{int\_in} and
\texttt{int\_out}, as described in
Section~\ref{sec:depfin-like}), eliminating these transfers.
The reduction is therefore not due to recomputation, but to the
\emph{elimination of intermediate data movement} between layers.
DRAM traffic is particularly informative for distinguishing
activation-dominant from weight-dominant workloads: for the former,
intermediate activations constitute the majority of off-chip traffic,
so fusion yields large savings; for the latter, weight transfers
dominate and the relative savings are smaller.
